\scanpage{122}\mspage{67}
L'idée me vient, après des essais \unsure{embrouillants|embrouillés}, que $\overline C$ pourrait être l'adjoint à \deleted{gauche}\inserted{droite}, et non à \deleted{droite}\inserted{gauche}, i.e. que les flèches cherchées pourraient aller en sens inverse
\[
 \begin{aligned}
 \mathrm{id}_{\unsure{A|\mathcal A}\Sigma^{-1}}&\xleftarrow{p}\unsure{\rho|p}\,\overline C\quad(=\overline{\rho_0 C}\deleted{\unsure{*}})\\
 \mathrm{id}_{A\Sigma_0^{-1}}&\xrightarrow{i}\overline C\rho\quad(=\overline{C\rho_0})
 \end{aligned}
\]
(j'inverse du coup les notations pour $p$ et $i$), donc on est amené si possible à chercher des flèches
\[
 \begin{aligned}
 \mathrm{id}_A&\mathrel{\deleted{\unsure{\rightarrow}}\!\!\xleftarrow{p}}\rho_0 C=C\\
 \mathrm{id}_A&\xrightarrow{i}C\rho_0=C,
 \end{aligned}
\]
donc des flèches fonctorielles en $X$
\[
 X\xrightarrow{i_X}C(X)\xrightarrow{p_X}X.
\]
On tient bel et bien un \deleted{$\unsure{p_X}$} $i_X$ fonctoriel --- tjrs le même ! Pour $p_X$, on n'en a pas plus qu'avant, mais se rappelant qu'il suffit de le trouver dans $A\Sigma^{-1}$ (et non plus $A\Sigma_0^{-1}$ !), il suffit d'inverser la flèche $i_X$, qui est en effet un \uline{iso} dans $A\Sigma^{-1}$ (mais non dans $A\Sigma_0^{-1}$). On posera donc
\[
 p_X=i_X^{-1}:C(X)\unsure{\text{---}|\longrightarrow}X\qquad\text{dans }A\Sigma^{-1}.
\]
Il faut prouver deux compatibilités seulement, pour les foncteurs $\overline C$ et $\rho$ entre $A\Sigma_0^{-1}$ et $A\Sigma^{-1}$ (et j'\deleted{\unsure{é}}écrirai $C$ au lieu de $\overline C$)\,: les deux composés
\diagram{\scriptsize\begin{tikzcd}[ampersand replacement=\&,column sep=small,row sep=large]
 {}\arrow[d,phantom,"{\Bigg\{}" description]\& \rho\arrow[r,"{\rho*i}"]\arrow[rr,bend left=28,dashed,"\mathrm{id}"]\&
 \rho C\rho\arrow[r,"{p*\rho}"]\&\rho\&\text{(relation \deleted{\unsure{pour les flèches|pour foncteurs}} à valeurs dans } A\Sigma^{-1}\text{)}\\
 {}\& C\arrow[r,"{i*C}"]\arrow[rr,bend right=23,dashed,"\mathrm{id}"']\&
 C\rho C\arrow[r,"{C*p}"]\& C\&\text{(relation pour foncteurs à valeurs dans } A\Sigma_0^{-1}\text{)}
\end{tikzcd}}
doivent être l'identité $\mathrm{id}_{\rho}$ resp $\mathrm{id}_C$.
% archive sheet number: 121
